%==========================================================
\section{Planeamento da investigação contextual}\label{sec:planeamento}

O planeamento seguiu o Capítulo~4 de \textit{Contextual Design}~\cite{beyer1998}, que propõe quatro passos: definir o foco, desenhar a investigação segundo o tipo de projeto, escolher a situação de entrevista adequada a cada tipo de tarefa e decidir quem entrevistar.

\subsection{Foco do projeto}

Os projetos tendem a definir-se pela solução que vão entregar, neste caso «um agente de viagens baseado em IA». Para saber quem entrevistar e o que observar, essa descrição tem de passar a descrever o \textbf{trabalho} que o sistema vai apoiar. Esse trabalho não é «usar uma aplicação de viagens». É todo o processo por que passa um jovem desde a vontade de viajar até voltar a casa com as contas acertadas.

\begin{foco}[Declaração de foco]
Como é que jovens com pouco dinheiro decidem para onde ir, descobrem e comparam opções, juntam as peças (transporte, dormida, atividades) dentro de um orçamento, combinam com quem vai e ajustam a viagem quando algo muda.
\end{foco}

A declaração usa linguagem simples de propósito: «juntar as peças» e não «agregação de serviços». Palavras simples levam o entrevistador a reparar num leque mais largo de comportamentos. A Tabela~\ref{tab:alargar} alarga o foco para lá da ferramenta.

\begin{table}[!htbp]
\centering\small
\begin{tabularx}{\textwidth}{@{}p{3.6cm}L@{}}
\toprule
\textbf{Pergunta} & \textbf{Aplicação ao \projeto} \\
\midrule
Que trabalho vamos apoiar? & Inspiração, escolha de destino, definição de datas flexíveis, pesquisa de transporte e alojamento, atividades, cálculo do custo total, reserva, divisão de despesas, alterações e imprevistos. \\
Quem está envolvido? & O viajante, os amigos que vão, quem organiza pelo grupo, os pais (que financiam ou autorizam), quem já foi ao destino, criadores de conteúdo e assistentes de IA já usados como ajudantes informais. \\
Onde acontece fisicamente? & No telemóvel, em momentos soltos; no portátil, em sessões longas com muitas abas; em grupos de mensagens; presencialmente na faculdade ou no café. \\
Em que contexto cultural? & Viajar barato vale como competência e motivo de orgulho. Há desconfiança de pacotes e de custos escondidos, pressão para não ficar de fora e o peso da opinião dos pais. \\
\bottomrule
\end{tabularx}
\caption{Perguntas para alargar o foco do projeto.}
\label{tab:alargar}
\end{table}

\paragraph{Trabalho análogo.} O livro sugere estudar atividades com a mesma estrutura noutros domínios. A mais próxima é o \textbf{agente de viagens de balcão}, que faz hoje o papel que o \projeto{} quer ocupar. Isso justificou entrevistar um agente (A0). Outras metáforas registadas para entrevistas futuras: montar um PC com orçamento fixo, organizar um jantar de grupo em que alguém adianta o dinheiro e fazer compras de supermercado com orçamento.

\subsection{Tipo de projeto}

O \projeto{} cruza dois dos pontos de partida descritos no livro:
\begin{itemize}
  \item \textbf{Produto conhecido.} Já existem comparadores de voos, sites de alojamento e de atividades, e o mercado já tem expectativas. É preciso observar jovens a usar estas ferramentas, registar o que atrapalha (taxas que só aparecem no fim, datas fixas, resultados separados por serviço) e procurar necessidades que ninguém reconheceu.
  \item \textbf{Nova tecnologia (agente de IA).} O análogo de um agente de IA é um agente humano, ou o amigo que «trata de tudo». Muitos jovens já pedem roteiros a assistentes de IA, e interessa ver o que fazem a seguir e onde deixam de confiar.
\end{itemize}

\begin{nota}[Princípio orientador]
«Todo o trabalho já está a ser feito de alguma forma.» É esse trabalho atual que estudamos, e estudá-lo não impede a inovação.
\end{nota}

\subsection{Situação de entrevista}

Planear uma viagem não é uma tarefa única e contínua. Mistura vários tipos de tarefa, e cada um exige uma técnica diferente (Tabela~\ref{tab:tarefas}).

\begin{table}[!htbp]
\centering\small
\begin{tabularx}{\textwidth}{@{}p{2.4cm}p{4.2cm}L@{}}
\toprule
\textbf{Tipo de tarefa} & \textbf{No planeamento de uma viagem} & \textbf{Técnica} \\
\midrule
Normal & Uma sessão de pesquisa no portátil. & Entrevista contextual com quem está a planear uma viagem real. \\
Intermitente & Ver um preço num intervalo, receber um link de um amigo. & Criar um rasto (capturas e notas durante 1--2 semanas) e fazer um relato retrospetivo. \\
Extremamente longa & O planeamento pode durar semanas. & \textit{Work walkthrough} de uma viagem já planeada, guiado pelos artefactos. \\
Interna (decisão) & «Vale a pena pagar mais 30\,€ por um voo direto?» & Criar o evento: pedir que planeie à nossa frente e interromper para propor interpretações. \\
Não interrompível & Discussão do grupo ao vivo ou por mensagens. & Não intervir; rever depois o histórico com o participante. \\
\bottomrule
\end{tabularx}
\caption{Tipos de tarefa e técnicas de entrevista.}
\label{tab:tarefas}
\end{table}

Cada entrevista segue o modelo mestre--aprendiz em quatro fases: \textbf{conversa convencional} (perfil e opinião geral, que não entra nos modelos sem caso concreto), \textbf{transição} (explicar as regras), \textbf{entrevista contextual} (observar, perguntar, manter tudo concreto) e \textbf{fecho} (resumir e deixar o participante corrigir).

\subsection{Quem entrevistar}

O livro recomenda 2--3 pessoas por papel e diversidade na forma de trabalhar, mais do que em segmentos demográficos. A Tabela~\ref{tab:papeis} mostra o plano e o que foi cumprido nesta entrega.

\begin{table}[!htbp]
\centering\small
\begin{tabularx}{\textwidth}{@{}p{3.6cm}Lcc@{}}
\toprule
\textbf{Papel} & \textbf{Descrição} & \textbf{Previsto} & \textbf{Feito} \\
\midrule
Organizador(a) & Planeia pelo grupo, adianta o dinheiro e cobra aos outros. & 3 & 1 (U0) \\
Quem vai atrás & Participa no grupo, vota e paga a sua parte, mas não pesquisa. & 2 & -- \\
Viajante solo & Decide tudo sozinho(a), com mais flexibilidade de datas. & 3 & -- \\
Caçador(a) de promoções & Parte do preço e não do destino. & 2 & (traços em U0) \\
Erasmus / interrail & Viagens com várias paragens a partir de outra cidade. & 2 & -- \\
Financiador(a) & Pai ou mãe que paga ou autoriza. & 1--2 & -- \\
Agente de viagens & Domínio da metáfora: faz hoje o que o \projeto{} quer fazer. & 1--2 & 1 (A0) \\
\midrule
\multicolumn{2}{@{}l}{\textbf{Total}} & \textbf{14--16} & \textbf{2} \\
\bottomrule
\end{tabularx}
\caption{Papéis a entrevistar: plano e execução.}
\label{tab:papeis}
\end{table}

As duas sessões realizadas cumprem o mínimo do enunciado e cobrem os dois lados do problema: quem planeia a viagem (U0) e quem faz hoje esse trabalho profissionalmente (A0). Os restantes papéis ficam para as próximas iterações.
